\section{A polynomially quantified mean-field prefactor window}
\label{sec:small-prefactor}

The bulk-subtraction proof depends on the prefactor $t$, not directly
on the exponent $\gamma$. This gives a separate finite-temperature
statement at mean-field scaling. It does not include the original
critical convention $t=1$.

\begin{theorem}[Small prefactor at every width]
\label{thm:small-meanfield-prefactor}
Assume the witness hypothesis of
Theorem~\ref{thm:unrestricted-above-one}, with parameters $d,L,C,K$
and $Q_*=[(K+2)/L]^{L/2}$. Define
\[
 A=\left(\frac{K+2}{L-1}\right)^{(L-1)/2},\qquad
 \eta=\min\left\{1,\frac1{4\sqrt{L(K+1)}},
                         \frac1{8\sqrt2\,LA}\right\},
\]
\[
 b^2=\max\{2(d+1)\log2+4\log8,\;16\log9+8\log(8L)\},
\qquad
 t_0=\min\left\{1,\frac\eta b,
                    \frac1{\sqrt{2\log(4/\varepsilon)}}\right\}.
\]
There is a sample bound polynomial in
$d,L,K,Q_*,\varepsilon^{-1},\log(2/\delta)$ such that, with probability
at least $1-\delta$, the unrestricted posterior has Gibbs error at most
$\varepsilon$, simultaneously for all $N\ge C$ and all
\[
 0<t\le t_0,\qquad \sigma^2=t^2/N.
\]
The inverse threshold $t_0^{-1}$ is polynomial in
$d,L,K,Q_*,\varepsilon^{-1}$. Neither bound depends on the realized
training sample or on its size through a circular condition.
\end{theorem}

\begin{proof}
Use the typical set $T$ and the perturbation estimate from
Section~\ref{sec:gaussian-distance}. For every width in the statement,
$tb\le\eta$, so the target shift has raw mass at least
$\frac12e^{-K/(2\sigma^2)}$. The posterior distance tail is therefore
\[
 Q_S\{d_\kappa(w,T)^2>K+1\}
 \le2e^{-1/(2\sigma^2)}
 =2e^{-N/(2t^2)}\le\varepsilon/2.
\]
Every remaining draw has exactly the same Boolean function as a
residual of energy at most $K+1$, with training margin greater than
$3/8$. The single width-independent generalization event underlying
\eqref{eq:distance-sample-bound} controls all of these residuals
simultaneously. Choosing its right-hand side at most $\varepsilon/2$
proves the claim. This sample choice is polynomial: for example, if
$G=\sqrt{(d+1)/d}[(K+1)/L]^{L/2}$, then a sufficiently large universal
constant times
\[
 [1+G+L+\varepsilon^{-1}+\log(4/\delta)]^{12}
\]
suffices. Finally $G\le\sqrt2Q_*$ and $A\le\sqrt eQ_*$ give all the
stated polynomial dependences. This proof requires no width split,
because the small-bulk condition has been imposed directly on $t$.
\end{proof}

\paragraph{A sample bound without an additional depth factor.}
The global-margin counting theorem gives a sharper choice of sample
size in Theorem~\ref{thm:small-meanfield-prefactor}. Put
\[
 \overline P=\max\left\{1,\frac8{3\sqrt d}
                  \left(\frac{K+1}{L}\right)^{L/2}\right\}.
\]
Rescale the residual readout by $8/3$. Its Boolean function now belongs
to $\mathcal H_{\overline P}$ from
Corollary~\ref{cor:all-depth-pf-counting}. Therefore it suffices to use
\[
 n\ge\frac2\varepsilon
   \left[p(\overline P)+\overline P^2\log(d+1)
                         +\log(2/\delta)\right],
\]
for a universal polynomial $p$. In particular, if $K\le L+B$ with
$B$ independent of the learner depth, then
\[
 \overline P\le\max\{1,(8/3)d^{-1/2}e^{(B+1)/2}\}.
\]
The required sample size has no remaining learner-depth dependence.
The sufficient prefactor $t_0$ still depends on $L$, and the original
critical choice $t=1$ is still excluded. Likewise, when this argument
is combined with the small-width fallback below, that finite-width
part retains its polynomial dependence on $L$.

\begin{corollary}[Polynomial product targets at mean-field scaling]
\label{cor:small-prefactor-product-targets}
Suppose $f$ has an arbitrary finite-width, depth-$D_0$ tanh
representation with unit global margin and product norm at most
$P\ge1$. There are universal polynomials $p,q$ and a universal
constant $c$ such that, whenever
\[
 L\ge D_0+c\log(2+d+D_0+P)
 \quad\hbox{or}\quad L\ge2D_0,
\]
sample size at least
\[
 p(d,D_0,P,L,\varepsilon^{-1},\log(2/\delta))
\]
gives Gibbs error at most $\varepsilon$ with probability at least
$1-\delta$, simultaneously for every actual width $N\ge C_L(f)$
and every mean-field scale
\[
 \sigma^2=t^2/N,
 \qquad 0<t\le
 q(d,D_0,P,L,\varepsilon^{-1},\log(2/\delta))^{-1}.
\]
For a bounded shallow target, or for a Boolean halfspace using
Theorem~\ref{thm:halfspace-unrestricted-all-depths}, the corresponding
statement holds at every $L\ge2$, with the same polynomial dependence
as in those positive theorems. The original critical prior $t=1$
is not asserted to satisfy this condition.
\end{corollary}

\begin{proof}
The width-compression and product-target constructions provide a
polynomial threshold $M$ and, at all $N\ge M$, unit-margin witnesses
with $K$ and $Q_*$ polynomial in $d,D_0,P,L$. Apply
Theorem~\ref{thm:small-meanfield-prefactor} to these widths. The physical
weight bound in that construction gives $C_L(f)\le M$.

For the remaining interval $C_L(f)\le N<M$, choose an actual
minimum-width bounded-weight witness and double its readout. It now
has unit margin and weights at most $2M$. Balanced cloning gives
unit-margin witnesses throughout this finite width interval with
energy polynomial in $d,L,M$. The finite-width theorem, with its
Gaussian-coordinate-box denominator, gives a polynomial sample bound
uniformly for all $t\le1$. Its constants can be bounded using $M$
and $2M$, independently of the actual minimum width. Union bound
this sample event and the large-width event, each with failure
probability $\delta/2$. Their thresholds are deterministic
polynomials of the displayed parameters. The same proof applies to
the bounded shallow and halfspace witnesses at all $L\ge2$.
\end{proof}
